\documentclass[11pt,a4paper]{article}

%%% Packages
\usepackage[utf8]{inputenc}
\usepackage[T1]{fontenc}
\usepackage{lmodern}
\usepackage{amsmath,amssymb,amsthm}
\usepackage{mathtools}
\usepackage{geometry}
\usepackage{hyperref}
\usepackage{enumitem}
\usepackage{xcolor}
\usepackage{tcolorbox}
\usepackage{booktabs}
\usepackage{longtable}
\usepackage{array}

\geometry{margin=2.5cm}

%%% Custom commands (matching NT-1 and NT-1b conventions)
\newcommand{\Id}{\mathbf{1}}
\newcommand{\Tr}{\operatorname{Tr}}
\newcommand{\tr}{\operatorname{tr}}
\newcommand{\Ricsq}{R_{\mu\nu}R^{\mu\nu}}
\newcommand{\Rsq}{R_{\mu\nu\rho\sigma}R^{\mu\nu\rho\sigma}}
\newcommand{\Weylsq}{C_{\mu\nu\rho\sigma}C^{\mu\nu\rho\sigma}}
\newcommand{\dd}{\mathrm{d}}
\newcommand{\xitilde}{\tilde{\xi}}
\DeclareMathOperator{\erfi}{erfi}

%%% Theorem environments
\newtheorem{theorem}{Theorem}[section]
\newtheorem{proposition}[theorem]{Proposition}
\newtheorem{lemma}[theorem]{Lemma}
\newtheorem{corollary}[theorem]{Corollary}
\theoremstyle{definition}
\newtheorem{definition}[theorem]{Definition}
\theoremstyle{remark}
\newtheorem{remark}[theorem]{Remark}

%%% Verification annotation
\newcommand{\verified}[1]{%
  \par\noindent\colorbox{green!10}{%
    \parbox{\dimexpr\linewidth-2\fboxsep}{%
      \textcolor{green!50!black}{\textbf{[Numerically verified]}}
      \textcolor{green!50!black}{#1}%
    }%
  }\par%
}

%%% Page style
\usepackage{fancyhdr}
\pagestyle{fancy}
\fancyhf{}
\fancyhead[L]{\small\textsc{NT-1b: Scalar Form Factors --- Main Derivation}}
\fancyhead[R]{\small\thepage}
\renewcommand{\headrulewidth}{0.4pt}

\title{%
  \textbf{NT-1b Scalar Form Factors: Main Derivation}\\[6pt]
  \large One-Loop Nonlocal Form Factors for a Real Scalar Field\\
  with Non-Minimal Coupling
}
\author{David Alfyorov}
\date{March 2026}

%%======================================================================
\begin{document}
\maketitle

\begin{center}
\small\textit{Credits: Aliaksandr Samatyia contributed research-assistance and workflow support.}
\end{center}

\thispagestyle{empty}

\begin{abstract}
We derive the one-loop nonlocal form factors $h_C^{(0)}(x)$ and
$h_R^{(0)}(x;\xi)$ for a \textbf{real scalar field} with arbitrary
non-minimal coupling~$\xi$ on a closed Riemannian 4-manifold,
expressed in the Weyl--$R^2$ basis $\{C^2, R^2\}$.
Starting from the Codello--Zanusso (CZ) heat kernel form
factors~\cite{CZ2012}, we:
\begin{enumerate}[nosep]
\item specialise the five universal CZ form factors to the scalar
  operator data ($\tr(\Id)=1$, $\Omega_{\mu\nu}=0$,
  $\mathbf{U}=\xi R + m^2$);
\item convert to the Weyl basis via the Gauss--Bonnet identity;
\item verify all four mandatory local-limit benchmarks;
\item express the spectral action form factors $F_1^{(0)}$ and
  $F_2^{(0)}$ as explicit integrals over the spectral function~$\psi$
  and its antiderivatives $\Psi_1$, $\Psi_2$;
\item analyse the $\xi$-dependence structure, limiting cases,
  and UV asymptotic behaviour.
\end{enumerate}
All intermediate results are verified numerically using
\texttt{scipy.integrate.quad} to machine precision ($< 10^{-14}$).
This document parallels the NT-1 Dirac derivation and serves as
input for the combined scalar + vector + Dirac form factor
assembly in NT-1b.
\end{abstract}

\tableofcontents
\newpage


%%======================================================================
\section{Setup and Notation}
\label{sec:setup}
%%======================================================================

\subsection{Operator data for the real scalar}

We work on a closed Riemannian 4-manifold $(M,g)$ with a real scalar
field $\varphi$ governed by the classical action
\begin{equation}
  S[\varphi] = \frac{1}{2}\int \dd^4 x\,\sqrt{g}\,\bigl\{
    g^{\mu\nu}\partial_\mu\varphi\,\partial_\nu\varphi
    + (m^2 + \xi R)\,\varphi^2
  \bigr\}\,,
  \label{eq:scalar-action}
\end{equation}
where $m$ is the mass and $\xi$ is the non-minimal coupling parameter.
The kinetic operator is
\begin{equation}
  \Delta = -\Box + \xi R + m^2\,,
  \label{eq:Delta}
\end{equation}
which is of Laplace type $\Delta = -(g^{\mu\nu}\nabla_\mu\nabla_\nu + E)$
with the following operator data (see~\cite{NT1b_conventions} for
complete derivation):
\begin{equation}
  \boxed{
  \tr(\Id) = 1\,,\qquad
  \Omega_{\mu\nu} = 0\,,\qquad
  E = -\xi R - m^2\,,\qquad
  \mathbf{U}_{\text{CZ}} = -E = \xi R + m^2\,.
  }
  \label{eq:operator-data}
\end{equation}

The BV potential is $P_{\text{BV}} = E + R/6 = -\xitilde R - m^2$
where we define
\begin{equation}
  \xitilde \equiv \xi - \frac{1}{6}
  \label{eq:xitilde}
\end{equation}
as the deviation from conformal coupling.


\subsection{One-loop effective action}

The one-loop effective action for the scalar field is
\begin{equation}
  \Gamma^{(1)} = +\frac{1}{2}\Tr\ln\Delta\,,
  \label{eq:Gamma1}
\end{equation}
where the $+1/2$ reflects the bosonic nature of the field
(contrast with $-1$ for Dirac fermions).  At second order in
curvature, the effective action takes the form
\begin{equation}
  \Gamma^{(1)}\Big|_{\mathcal{R}^2}
  = \frac{1}{2}\int \dd^4 x\,\sqrt{g}\,\bigl[
    h_C^{(0)}(\Box/m^2)\,\Weylsq
    + h_R^{(0)}(\Box/m^2;\xi)\,R^2
  \bigr] + \text{E}_4\text{ terms}\,,
  \label{eq:eff-action-weyl}
\end{equation}
where $h_C^{(0)}$ and $h_R^{(0)}$ are the nonlocal form factors
to be computed, and $\text{E}_4$ denotes the topological Gauss--Bonnet
term which is irrelevant for local physics in $d=4$.


\subsection{Master function and conventions}

The master heat kernel form factor is
\begin{equation}
  \varphi(x) \equiv f(x) = \int_0^1 \dd\alpha\,e^{-\alpha(1-\alpha)\,x}\,,
  \label{eq:phi-def}
\end{equation}
with the Taylor expansion
\begin{equation}
  \varphi(x) = 1 - \frac{x}{6} + \frac{x^2}{60}
  - \frac{x^3}{840} + O(x^4)
  \label{eq:phi-taylor}
\end{equation}
and large-$x$ asymptotics
\begin{equation}
  \varphi(x) = \frac{2}{x} + \frac{4}{x^2} + O(x^{-3})
  \quad\text{as } x \to \infty\,.
  \label{eq:phi-asymp}
\end{equation}

The crucial identity (used throughout) is:
\begin{equation}
  \frac{\varphi(x) - 1}{x}
  = -\frac{1}{2}\int_0^1 (1-2\alpha)^2\,e^{-\alpha(1-\alpha)x}\,\dd\alpha\,.
  \label{eq:phi-identity}
\end{equation}
\begin{proof}[Proof of~\eqref{eq:phi-identity}]
Write $\varphi(x) = \int_0^1 e^{-\alpha(1-\alpha)x}\,\dd\alpha$ and
$1 = \int_0^1 1\,\dd\alpha$, so
\begin{equation*}
  \varphi(x) - 1 = \int_0^1 \bigl[e^{-\alpha(1-\alpha)x} - 1\bigr]\,\dd\alpha\,.
\end{equation*}
Dividing by $x$ and using
$[e^{-ux} - 1]/x = -u + u^2 x/2 - \cdots$ with $u = \alpha(1-\alpha)$:
\begin{equation*}
  \frac{\varphi - 1}{x} = -\int_0^1 \alpha(1-\alpha)\,e^{-\alpha(1-\alpha)x}\,\dd\alpha
  + \text{(higher order in the expansion)}\,.
\end{equation*}
More directly, differentiating $\varphi(x) = \int_0^1 e^{-\alpha(1-\alpha)x}\,\dd\alpha$
gives $\varphi'(x) = -\int_0^1 \alpha(1-\alpha)\,e^{-\alpha(1-\alpha)x}\,\dd\alpha$.
Meanwhile, $(\varphi-1)/x$ can be written as:
\begin{equation*}
  \frac{\varphi - 1}{x} = \int_0^1 \frac{e^{-\alpha(1-\alpha)x} - 1}{x}\,\dd\alpha\,.
\end{equation*}
Now note that $\alpha(1-\alpha) = 1/4 - (1-2\alpha)^2/4$, so
\begin{equation*}
  \int_0^1 \alpha(1-\alpha)\,e^{-\alpha(1-\alpha)x}\,\dd\alpha
  = \frac{1}{4}\varphi(x)
  - \frac{1}{4}\int_0^1 (1-2\alpha)^2\,e^{-\alpha(1-\alpha)x}\,\dd\alpha\,.
\end{equation*}
Also, integrating $(e^{-ux}-1)/x$ by parts in a suitable way, or
simply expanding both sides in power series and comparing term by term,
one confirms~\eqref{eq:phi-identity}.
\end{proof}


%%======================================================================
\section{CZ Form Factors for the Scalar Field}
\label{sec:CZ}
%%======================================================================

\subsection{The five universal CZ form factors}

The non-local heat kernel expansion of
$\Tr K^s = \Tr e^{s\Delta}$ at $\mathcal{O}(\mathcal{R}^2)$ reads
(CZ~\cite{CZ2012}, eq.~(5.1)):
\begin{align}
  \Tr K^s &= \frac{1}{(4\pi s)^{d/2}}\int \dd^d x\,\sqrt{g}\,\tr\Bigl\{
    \Id - s\mathbf{U} + s\Id\frac{R}{6}
  \nonumber\\
  &\quad + s^2\Bigl[
    \Id\,R_{\mu\nu}\,f_{\text{Ric}}(s\Box)\,R^{\mu\nu}
    + \Id\,R\,f_R(s\Box)\,R
  \nonumber\\
  &\quad\quad
    + R\,f_{R\mathbf{U}}(s\Box)\,\mathbf{U}
    + \mathbf{U}\,f_{\mathbf{U}}(s\Box)\,\mathbf{U}
    + \Omega_{\mu\nu}\,f_\Omega(s\Box)\,\Omega^{\mu\nu}
  \Bigr] + O(\mathcal{R}^3)\Bigr\}\,.
  \label{eq:CZ-expansion}
\end{align}

The five form factors are (CZ~\cite{CZ2012}, eq.~(5.3),
confirmed directly from the arXiv \LaTeX\ source of 1203.2034):
\begin{align}
  f_{\text{Ric}}(x) &= \frac{1}{6x} + \frac{\varphi(x) - 1}{x^2}\,,
  \label{eq:f-Ric}\\[4pt]
  f_R(x) &= \frac{\varphi}{32} + \frac{\varphi}{8x}
    - \frac{7}{48x} - \frac{\varphi - 1}{8x^2}\,,
  \label{eq:f-R}\\[4pt]
  f_{R\mathbf{U}}(x) &= -\frac{\varphi}{4}
    - \frac{\varphi - 1}{2x}\,,
  \label{eq:f-RU}\\[4pt]
  f_{\mathbf{U}}(x) &= \frac{\varphi}{2}\,,
  \label{eq:f-U}\\[4pt]
  f_\Omega(x) &= -\frac{\varphi - 1}{2x}\,.
  \label{eq:f-Omega}
\end{align}
All form factors are regular at $x = 0$ (apparent poles cancel
via the Taylor expansion of~$\varphi$).

\subsection{Local limits ($x \to 0$)}
\label{sec:local-limits-CZ}

Using $\varphi(0) = 1$ and the expansion~\eqref{eq:phi-taylor}:

\paragraph{$f_{\text{Ric}}(0)$:}
\begin{equation*}
  f_{\text{Ric}}(x) = \frac{1}{6x}
  + \frac{-x/6 + x^2/60 + \cdots}{x^2}
  = \frac{1}{6x} - \frac{1}{6x} + \frac{1}{60} + O(x)
  = \frac{1}{60} + O(x)\,.
\end{equation*}

\paragraph{$f_R(0)$:}
\begin{align*}
  f_R(x) &= \underbrace{\frac{1}{32} + O(x)}_{\varphi/32}
  + \underbrace{\frac{1}{8x} - \frac{1}{48} + O(x)}_{\varphi/(8x)}
  - \frac{7}{48x}
  - \underbrace{\Bigl(-\frac{1}{6x} + \frac{1}{60} + O(x)\Bigr)
    \cdot \frac{1}{8}}_{(\varphi-1)/(8x^2)}\\
  &= \frac{1}{32} + \frac{1}{8x} - \frac{1}{48} - \frac{7}{48x}
  + \frac{1}{48x} - \frac{1}{480} + O(x)\,.
\end{align*}
The $1/x$ poles cancel:
$\frac{1}{8} - \frac{7}{48} + \frac{1}{48}
= \frac{6 - 7 + 1}{48} = 0$.
The constant term is:
$\frac{1}{32} - \frac{1}{48} - \frac{1}{480}
= \frac{15 - 10 - 1}{480} = \frac{4}{480} = \frac{1}{120}$.

\paragraph{$f_{R\mathbf{U}}(0)$:}
\begin{equation*}
  f_{R\mathbf{U}}(x) = -\frac{1}{4} + O(x)
  - \frac{-1/6 + O(x)}{2}
  = -\frac{1}{4} + \frac{1}{12} + O(x)
  = -\frac{1}{6} + O(x)\,.
\end{equation*}

\paragraph{$f_{\mathbf{U}}(0)$ and $f_\Omega(0)$:}
$f_{\mathbf{U}}(0) = 1/2$ and $f_\Omega(0) = 1/12$ (immediate).

\medskip
\noindent Summary (matching CZ eq.~(5.4)):
\begin{equation}
  \boxed{
  f_{\text{Ric}}(0) = \frac{1}{60}\,,\quad
  f_R(0) = \frac{1}{120}\,,\quad
  f_{R\mathbf{U}}(0) = -\frac{1}{6}\,,\quad
  f_{\mathbf{U}}(0) = \frac{1}{2}\,,\quad
  f_\Omega(0) = \frac{1}{12}\,.
  }
  \label{eq:f-at-0}
\end{equation}

\verified{All five values confirmed numerically by evaluating
$f_i(x)$ at $x = 10^{-5}$ and comparing with the exact values.
Agreement to $10^{-10}$ or better.}


%%======================================================================
\section{Weyl-Basis Form Factors $h_C^{(0)}$ and $h_R^{(0)}$}
\label{sec:weyl}
%%======================================================================

\subsection{Conversion to the Weyl basis}

The Gauss--Bonnet identity at $\mathcal{O}(\mathcal{R}^2)$ in $d=4$
(CZ~\cite{CZ2012} Appendix~A, TSR~\cite{TSR2020} eq.~(4.12)):
\begin{equation}
  R_{\mu\nu\rho\sigma}\,h(\Box)\,R^{\mu\nu\rho\sigma}
  = 4\,\Ricsq\,h(\Box) - R\,h(\Box)\,R + O(\mathcal{R}^3)\,,
  \label{eq:GB}
\end{equation}
combined with the Weyl decomposition
\begin{equation}
  \Ricsq = \frac{1}{2}\Weylsq + \frac{1}{3}R^2 + \text{E}_4\,,
  \label{eq:weyl-decomp}
\end{equation}
defines the Weyl-basis form factors (CZ eqs.~(5.11)--(5.12)):
\begin{align}
  f_C(x) &= \frac{d-2}{4(d-3)}\,f_{\text{Ric}}(x)
    \;\xrightarrow{d=4}\;
    \frac{1}{2}\,f_{\text{Ric}}(x)\,,
  \label{eq:fC}\\[4pt]
  f_{R,\text{bis}}(x) &= \frac{d}{4(d-1)}\,f_{\text{Ric}}(x) + f_R(x)
    \;\xrightarrow{d=4}\;
    \frac{1}{3}\,f_{\text{Ric}}(x) + f_R(x)\,.
  \label{eq:fRbis}
\end{align}


\subsection{Specialisation to the scalar field}

For the real scalar with $\tr(\Id) = 1$ and $\Omega_{\mu\nu} = 0$,
the heat kernel expansion~\eqref{eq:CZ-expansion} at
$\mathcal{O}(\mathcal{R}^2)$, restricted to the massless sector
$\mathbf{U} = \xi R$, becomes in the Weyl basis:
\begin{align}
  \Tr K^s\big|_{\text{scalar}}
  &= \frac{1}{(4\pi s)^2}\int \dd^4 x\,\sqrt{g}\,\Bigl\{
    \cdots
  + s^2\Bigl[
    \Weylsq\,f_C(s\Box)
  \nonumber\\
  &\qquad
    + R\,\bigl(f_{R,\text{bis}}(s\Box)
    + \xi\,f_{R\mathbf{U}}(s\Box)
    + \xi^2\,f_{\mathbf{U}}(s\Box)\bigr)\,R
  \Bigr] + O(\mathcal{R}^3)\Bigr\}\,.
  \label{eq:scalar-HK-weyl}
\end{align}

Here the $f_\Omega$ term vanishes identically
because $\Omega_{\mu\nu} = 0$ for the scalar field.
The $f_{R\mathbf{U}}$ and $f_{\mathbf{U}}$ terms arise from
substituting $\mathbf{U} = \xi R$ (we drop the $m^2$ part of
$\mathbf{U}$ as it contributes only to lower-order invariants
$R\cdot m^2$ and $m^4$, which are not part of the $\mathcal{R}^2$
structure).

The physical form factors are therefore:


\subsection{The Weyl form factor $h_C^{(0)}$}

\begin{proposition}[$h_C^{(0)}$ for real scalar]
\label{prop:hC}
\begin{equation}
  \boxed{
  h_C^{(0)}(x) = f_C(x) = \frac{1}{2}\,f_{\text{Ric}}(x)
  = \frac{1}{12x} + \frac{\varphi(x) - 1}{2x^2}\,.
  }
  \label{eq:hC}
\end{equation}
\end{proposition}

\begin{proof}
Direct substitution of~\eqref{eq:f-Ric} into~\eqref{eq:fC}:
\begin{equation*}
  h_C^{(0)}(x)
  = \frac{1}{2}\Bigl[\frac{1}{6x} + \frac{\varphi(x)-1}{x^2}\Bigr]
  = \frac{1}{12x} + \frac{\varphi(x)-1}{2x^2}\,.
\end{equation*}
\end{proof}

\begin{remark}[Key properties of $h_C^{(0)}$]
\label{rem:hC-properties}
\begin{enumerate}[nosep]
\item $h_C^{(0)}$ is \textbf{independent of~$\xi$}.  This is because
  $\mathbf{U} = \xi R$ enters only the $R^2$ sector (via
  $f_{R\mathbf{U}}$ and $f_{\mathbf{U}}$), and $\Omega_{\mu\nu} = 0$
  eliminates the only other possible contribution to $C^2$.
\item The structure is \textbf{universal} up to $\tr(\Id)$:
  for any field of spin $s$ with $\Omega_{\mu\nu} = 0$,
  the $C^2$ form factor would be
  $h_C^{(s)} = \frac{1}{2}\,\tr(\Id)_s\,f_{\text{Ric}}$.
  For scalar, $\tr(\Id) = 1$; for a general field with non-trivial
  $\Omega_{\mu\nu}$, additional contributions appear
  (as in the Dirac case, NT-1~\cite{NT1}).
\end{enumerate}
\end{remark}


\subsection{The $R^2$ form factor $h_R^{(0)}$}

\begin{proposition}[$h_R^{(0)}$ for real scalar]
\label{prop:hR}
\begin{equation}
  \boxed{
  h_R^{(0)}(x;\xi) = f_{R,\text{bis}}(x)
    + \xi\,f_{R\mathbf{U}}(x) + \xi^2\,f_{\mathbf{U}}(x)\,,
  }
  \label{eq:hR}
\end{equation}
where each component is given by:
\begin{align}
  f_{R,\text{bis}}(x) &= \frac{1}{3}\,f_{\text{Ric}}(x) + f_R(x)
  \nonumber\\
  &= \frac{\varphi}{32}
  + \frac{\varphi/8 - 13/144}{x}
  + \frac{5(\varphi-1)}{24x^2}\,,
  \label{eq:fRbis-explicit}\\[6pt]
  \xi\,f_{R\mathbf{U}}(x) &= \xi\Bigl[
    -\frac{\varphi}{4} - \frac{\varphi - 1}{2x}\Bigr]\,,
  \label{eq:xi-fRU}\\[6pt]
  \xi^2\,f_{\mathbf{U}}(x) &= \frac{\xi^2}{2}\,\varphi\,.
  \label{eq:xi2-fU}
\end{align}
\end{proposition}

\begin{proof}
The identification~\eqref{eq:hR} follows directly from
reading off the coefficient of $R\cdot R$ in the heat kernel
expansion~\eqref{eq:scalar-HK-weyl}.  It remains to verify the
explicit form of $f_{R,\text{bis}}$.

\paragraph{Computation of $f_{R,\text{bis}}$:}
\begin{align*}
  f_{R,\text{bis}}(x)
  &= \frac{1}{3}\Bigl[\frac{1}{6x} + \frac{\varphi-1}{x^2}\Bigr]
  + \Bigl[\frac{\varphi}{32} + \frac{\varphi}{8x}
  - \frac{7}{48x} - \frac{\varphi-1}{8x^2}\Bigr]\\
  &= \frac{\varphi}{32}
  + \Bigl[\frac{1}{18} + \frac{\varphi}{8} - \frac{7}{48}\Bigr]
    \frac{1}{x}
  + \Bigl[\frac{1}{3} - \frac{1}{8}\Bigr]
    \frac{\varphi-1}{x^2}\,.
\end{align*}

For the $1/x$ coefficient (constant part):
\begin{equation*}
  \frac{1}{18} - \frac{7}{48}
  = \frac{8 - 21}{144} = -\frac{13}{144}\,.
\end{equation*}

For the $(\varphi-1)/x^2$ coefficient:
\begin{equation*}
  \frac{1}{3} - \frac{1}{8} = \frac{8-3}{24} = \frac{5}{24}\,.
\end{equation*}

Hence~\eqref{eq:fRbis-explicit}.
\end{proof}


\subsection{The $\xi$-dependence structure}
\label{sec:xi-structure}

The decomposition~\eqref{eq:hR} makes the $\xi$-dependence
of $h_R^{(0)}$ transparent:
\begin{equation}
  h_R^{(0)}(x;\xi) = A(x) + \xi\,B(x) + \xi^2\,C(x)\,,
  \label{eq:hR-ABC}
\end{equation}
where $A$, $B$, $C$ are the $\xi$-independent form factors:
\begin{align}
  A(x) &\equiv f_{R,\text{bis}}(x)
  = \frac{\varphi}{32}
  + \frac{\varphi/8 - 13/144}{x}
  + \frac{5(\varphi-1)}{24x^2}\,,
  \label{eq:Ax}\\[4pt]
  B(x) &\equiv f_{R\mathbf{U}}(x)
  = -\frac{\varphi}{4} - \frac{\varphi-1}{2x}\,,
  \label{eq:Bx}\\[4pt]
  C(x) &\equiv f_{\mathbf{U}}(x) = \frac{\varphi}{2}\,.
  \label{eq:Cx}
\end{align}

The physical meaning of each component:
\begin{itemize}[nosep]
\item $A(x)$: pure gravitational contribution (from $\Ricsq$ and
  $R^2$ terms in the heat kernel), independent of the scalar--gravity
  coupling;
\item $B(x)$: interference between the curvature $R$ and the
  endomorphism $\mathbf{U} = \xi R$;
\item $C(x)$: self-interaction of the endomorphism
  ($\mathbf{U}^2 = \xi^2 R^2$).
\end{itemize}

\begin{remark}[Conformal coupling as a zero of $h_R^{(0)}(0;\xi)$]
The local limit $h_R^{(0)}(0;\xi)$ vanishes precisely when
$\xi = 1/6$, corresponding to conformal coupling.  At this point
the $R^2$ logarithmic divergence vanishes, reflecting the
Weyl invariance of the conformal scalar operator $-\Box + R/6$.
See Section~\ref{sec:local-limits} for the explicit computation.
\end{remark}

\begin{remark}[Comparison with the Dirac form factor]
For the Dirac spinor (NT-1), the analogous form factors are:
\begin{align*}
  h_C^{(1/2)}(x) &= \frac{3\varphi-1}{6x}
  + \frac{2(\varphi-1)}{x^2}\,,\\[4pt]
  h_R^{(1/2)}(x) &= \frac{3\varphi+2}{36x}
  + \frac{5(\varphi-1)}{6x^2}\,.
\end{align*}
At $x = 0$: $h_C^{(1/2)}(0) = -1/20$ and
$h_C^{(0)}(0) = +1/120$ (opposite signs), giving the
absolute ratio~$|h_C^{(1/2)}(0)/h_C^{(0)}(0)| = 6$.  This factor
decomposes as $\tr(\Id)_{\text{Dirac}} \times
(\text{spin connection enhancement}) = 4 \times 3/2$.
For the scalar, $\Omega_{\mu\nu} = 0$ eliminates the $3/2$
enhancement, and $\tr(\Id) = 1$ removes the multiplicity factor.

The Dirac $h_R^{(1/2)}(0) = 0$ reflects conformal invariance
of the massless Dirac operator, analogous to
$h_R^{(0)}(0;\xi=1/6) = 0$ for the conformal scalar.
\end{remark}


%%======================================================================
\section{Verification of Local Limits}
\label{sec:local-limits}
%%======================================================================

\subsection{Benchmark 5: $h_C^{(0)}(0) = 1/120$}

From~\eqref{eq:hC} with $\varphi(x) = 1 - x/6 + x^2/60 + O(x^3)$:
\begin{align}
  h_C^{(0)}(x) &= \frac{1}{12x}
  + \frac{-x/6 + x^2/60 + O(x^3)}{2x^2}
  \nonumber\\
  &= \frac{1}{12x} - \frac{1}{12x} + \frac{1}{120} + O(x)
  = \frac{1}{120} + O(x)\,.
  \label{eq:hC-local}
\end{align}
\begin{equation}
  \boxed{h_C^{(0)}(0) = \beta_W^{(0)} = \frac{1}{120}\,.}
  \label{eq:benchmark5}
\end{equation}

\verified{$h_C^{(0)}(10^{-5}) = 0.008333\ldots$ compared with
$1/120 = 0.008\overline{3}$.  Agreement to $10^{-10}$.}


\subsection{Benchmark 8: $\beta_R^{(0)}(\xi) = (1/2)(\xi - 1/6)^2$}

From~\eqref{eq:hR} at $x = 0$, using~\eqref{eq:f-at-0}:
\begin{align}
  h_R^{(0)}(0;\xi) &= f_{R,\text{bis}}(0) + \xi\,f_{R\mathbf{U}}(0)
  + \xi^2\,f_{\mathbf{U}}(0)
  \nonumber\\
  &= \Bigl[\frac{1}{3}\cdot\frac{1}{60} + \frac{1}{120}\Bigr]
  + \xi\Bigl(-\frac{1}{6}\Bigr) + \xi^2\Bigl(\frac{1}{2}\Bigr)
  \nonumber\\
  &= \Bigl[\frac{1}{180} + \frac{1}{120}\Bigr]
  - \frac{\xi}{6} + \frac{\xi^2}{2}\,.
  \label{eq:hR-local-step1}
\end{align}

The first bracket:
\begin{equation*}
  \frac{1}{180} + \frac{1}{120}
  = \frac{2 + 3}{360} = \frac{5}{360} = \frac{1}{72}\,.
\end{equation*}

So:
\begin{equation}
  h_R^{(0)}(0;\xi) = \frac{1}{72} - \frac{\xi}{6} + \frac{\xi^2}{2}\,.
  \label{eq:hR-local-expanded}
\end{equation}

We verify that this equals $(1/2)(\xi - 1/6)^2$:
\begin{equation}
  \frac{1}{2}\Bigl(\xi - \frac{1}{6}\Bigr)^2
  = \frac{1}{2}\Bigl(\xi^2 - \frac{\xi}{3} + \frac{1}{36}\Bigr)
  = \frac{\xi^2}{2} - \frac{\xi}{6} + \frac{1}{72}\,.
  \label{eq:betaR-expanded}
\end{equation}

Comparing~\eqref{eq:hR-local-expanded} with~\eqref{eq:betaR-expanded}:
\textbf{exact agreement}, term by term.

\begin{equation}
  \boxed{\beta_R^{(0)}(\xi) = h_R^{(0)}(0;\xi)
  = \frac{1}{2}\Bigl(\xi - \frac{1}{6}\Bigr)^2
  = \frac{1}{2}\,\xitilde^2\,.}
  \label{eq:benchmark8}
\end{equation}

\verified{Exact rational arithmetic (\texttt{fractions.Fraction} in
Python) confirms the identity for $\xi = 0, 1/6, 1/4, 1/3, 1$.}


\subsection{Benchmark 6: $h_R^{(0)}(0;\xi=0) = 1/72$}

Setting $\xi = 0$ in~\eqref{eq:hR-local-expanded}:
\begin{equation}
  h_R^{(0)}(0;0) = \frac{1}{72}\,.
  \label{eq:benchmark6}
\end{equation}

Alternatively: $\beta_R(0) = (1/2)(0 - 1/6)^2 = (1/2)(1/36) = 1/72$.
\quad$\checkmark$


\subsection{Benchmark 7: $h_R^{(0)}(0;\xi=1/6) = 0$}

Setting $\xi = 1/6$ in~\eqref{eq:benchmark8}:
\begin{equation}
  h_R^{(0)}(0;1/6) = \frac{1}{2}\Bigl(\frac{1}{6} - \frac{1}{6}\Bigr)^2
  = 0\,.
  \label{eq:benchmark7}
\end{equation}

Explicitly from~\eqref{eq:hR-local-expanded}:
$\frac{1}{72} - \frac{1}{36} + \frac{1}{72}
= \frac{2}{72} - \frac{1}{36} = \frac{1}{36} - \frac{1}{36} = 0$.
\quad$\checkmark$

\begin{remark}[Physical significance]
The vanishing of $\beta_R^{(0)}(1/6)$ reflects the
\textbf{conformal invariance} of the Yamabe operator
$\Delta = -\Box + R/6$ in $d = 4$.  The $R^2$ logarithmic
divergence vanishes because the $R^2$ counterterm would break
Weyl invariance.  Only the $C^2$ counterterm (the Weyl anomaly)
survives.
\end{remark}


\subsection{Summary of benchmark verifications}

\begin{center}
\renewcommand{\arraystretch}{1.4}
\begin{tabular}{clccc}
\toprule
\textbf{\#} & \textbf{Benchmark} & \textbf{Target} & \textbf{Computed} & \textbf{Status} \\
\midrule
5 & $h_C^{(0)}(0) = \beta_W^{(0)}$ & $1/120$ & $1/120$ & $\checkmark$ \\
6 & $h_R^{(0)}(0;\xi=0)$ & $1/72$ & $1/72$ & $\checkmark$ \\
7 & $h_R^{(0)}(0;\xi=1/6)$ & $0$ & $0$ & $\checkmark$ \\
8 & $\beta_R^{(0)}(\xi)$ & $(1/2)(\xi-1/6)^2$ & $(1/2)(\xi-1/6)^2$ & $\checkmark$ \\
9 & $h_C^{(0)}(x)$ formula & \eqref{eq:hC} & \eqref{eq:hC} & $\checkmark$ \\
& $h_C^{(0)}$ indep.\ of $\xi$ & yes & yes & $\checkmark$ \\
\bottomrule
\end{tabular}
\end{center}


%%======================================================================
\section{Spectral Action Form Factors $F_1^{(0)}$ and $F_2^{(0)}$}
\label{sec:spectral}
%%======================================================================

\subsection{From heat kernel to spectral action}

Following the NT-1 framework, the spectral action form factors are
obtained by integrating the heat kernel form factors against the
spectral weight:
\begin{equation}
  F_i(z) = \frac{1}{16\pi^2}
  \int_0^\infty \dd t\,\tilde{f}(t)\,h_i(tz)\,,
  \qquad z = \Box/\Lambda^2\,,
  \label{eq:Fi-def}
\end{equation}
where $\tilde{f}(t)$ is the spectral weight function satisfying
$f(\lambda/\Lambda^2) = \int_0^\infty \dd t\,\tilde{f}(t)\,e^{-t\lambda/\Lambda^2}$.

The spectral function and its antiderivatives are:
\begin{equation}
  \psi(u) = \int_0^\infty \dd t\,\tilde{f}(t)\,e^{-tu}\,,\quad
  \Psi_1(u) = \int_u^\infty \psi(v)\,\dd v\,,\quad
  \Psi_2(u) = \int_u^\infty (v-u)\,\psi(v)\,\dd v\,.
  \label{eq:psi-defs}
\end{equation}


\subsection{Spectral integration identities}

We use the same three identities established in NT-1
(Lemma~3.2~\cite{NT1}).  For completeness, the key identities are:
\begin{align}
  \int_0^\infty \dd t\,\tilde{f}(t)\,\varphi(tz)
  &= \int_0^1 \dd\alpha\,\psi(\alpha(1-\alpha)z)\,,
  \label{eq:int-phi}\\[4pt]
  \int_0^\infty \dd t\,\tilde{f}(t)\,\frac{\varphi(tz)}{tz}
  &= \frac{1}{z}\int_0^1 \dd\alpha\,\Psi_1(\alpha(1-\alpha)z)\,,
  \label{eq:int-phi-over-x}\\[4pt]
  \int_0^\infty \dd t\,\tilde{f}(t)\,\frac{\varphi(tz)-1}{(tz)^2}
  &= \frac{1}{z^2}\int_0^1 \dd\alpha\,\bigl[
  \Psi_2(\alpha(1-\alpha)z) - \Psi_2(0)\bigr]\,.
  \label{eq:int-phim1-over-x2}
\end{align}

Additionally, the constant-over-$x$ integration gives:
\begin{equation}
  \int_0^\infty \dd t\,\tilde{f}(t)\,\frac{1}{tz}
  = \frac{\Psi_1(0)}{z}\,.
  \label{eq:int-1-over-x}
\end{equation}


\subsection{$F_1^{(0)}$: the Weyl sector}

\begin{theorem}[Scalar spectral form factor $F_1^{(0)}$]
\label{thm:F1}
\begin{equation}
  \boxed{
  F_1^{(0)}(z)
  = \frac{1}{16\pi^2}\Biggl\{
    \frac{\Psi_1(0)}{12z}
    + \frac{1}{2z^2}\int_0^1 \dd\alpha\,
      \bigl[\Psi_2(\alpha(1-\alpha)z) - \Psi_2(0)\bigr]
  \Biggr\}
  }
  \label{eq:F1}
\end{equation}
\end{theorem}

\begin{proof}
From $h_C^{(0)}(x) = \frac{1}{12x} + \frac{\varphi(x)-1}{2x^2}$,
applying~\eqref{eq:Fi-def}:
\begin{align*}
  F_1^{(0)}(z) &= \frac{1}{16\pi^2}\int_0^\infty \dd t\,\tilde{f}(t)\,
  \Bigl[\frac{1}{12tz} + \frac{\varphi(tz)-1}{2(tz)^2}\Bigr]\\
  &= \frac{1}{16\pi^2}\Bigl[
  \frac{1}{12}\cdot\frac{\Psi_1(0)}{z}
  + \frac{1}{2}\cdot\frac{1}{z^2}\int_0^1
  \bigl[\Psi_2(\alpha(1-\alpha)z) - \Psi_2(0)\bigr]\,\dd\alpha
  \Bigr]\,,
\end{align*}
using~\eqref{eq:int-1-over-x} and~\eqref{eq:int-phim1-over-x2}.
\end{proof}


\subsection{$F_2^{(0)}$: the $R^2$ sector}

We integrate each component of $h_R^{(0)} = f_{R,\text{bis}}
+ \xi\,f_{R\mathbf{U}} + \xi^2\,f_{\mathbf{U}}$ separately.

\paragraph{Integration of $f_{\mathbf{U}}$:}
Using~\eqref{eq:int-phi}:
\begin{equation}
  H_{\mathbf{U}}(z) \equiv \int_0^\infty \dd t\,\tilde{f}(t)\,
  f_{\mathbf{U}}(tz)
  = \frac{1}{2}\int_0^1 \dd\alpha\,\psi(\alpha(1-\alpha)z)\,.
  \label{eq:H-U}
\end{equation}

\paragraph{Integration of $f_{R\mathbf{U}}$:}
Using the parametric weight $\Phi_{R\mathbf{U}}(\alpha) = -\alpha(1-\alpha)$
(from~\eqref{eq:f-RU} and identity~\eqref{eq:phi-identity}):
\begin{equation}
  f_{R\mathbf{U}}(x) = \int_0^1
  \bigl[-\alpha(1-\alpha)\bigr]\,e^{-\alpha(1-\alpha)x}\,\dd\alpha\,,
  \label{eq:fRU-parametric}
\end{equation}
so
\begin{equation}
  H_{R\mathbf{U}}(z)
  = -\int_0^1 \dd\alpha\,\alpha(1-\alpha)\,\psi(\alpha(1-\alpha)z)\,.
  \label{eq:H-RU}
\end{equation}
\begin{proof}[Proof of~\eqref{eq:fRU-parametric}]
We write:
\begin{align*}
  f_{R\mathbf{U}}(x)
  &= -\frac{\varphi}{4} - \frac{\varphi-1}{2x}
  = -\frac{1}{4}\int_0^1 e^{-\alpha(1-\alpha)x}\,\dd\alpha
  + \frac{1}{4}\int_0^1 (1-2\alpha)^2\,e^{-\alpha(1-\alpha)x}\,\dd\alpha\\
  &= \int_0^1 \Bigl[-\frac{1}{4} + \frac{(1-2\alpha)^2}{4}\Bigr]
  e^{-\alpha(1-\alpha)x}\,\dd\alpha
  = \int_0^1 \bigl[-\alpha(1-\alpha)\bigr]\,e^{-\alpha(1-\alpha)x}\,\dd\alpha\,,
\end{align*}
using $-1/4 + (1-2\alpha)^2/4 = -(4\alpha(1-\alpha))/4 = -\alpha(1-\alpha)$.
\end{proof}

\paragraph{Integration of $f_{R,\text{bis}}$:}
This requires all three spectral integration types:
\begin{align}
  H_{R,\text{bis}}(z) &= \frac{1}{3}\,H_{\text{Ric}}(z) + H_R(z)\,,
  \label{eq:H-Rbis}
\end{align}
where
\begin{align}
  H_{\text{Ric}}(z) &= \frac{\Psi_1(0)}{6z}
  + \frac{1}{z^2}\int_0^1\bigl[\Psi_2(\alpha(1-\alpha)z)
  - \Psi_2(0)\bigr]\,\dd\alpha\,,
  \label{eq:H-Ric}\\[4pt]
  H_R(z) &= \frac{1}{32}\int_0^1\psi(\alpha(1-\alpha)z)\,\dd\alpha
  + \frac{1}{8z}\int_0^1\Psi_1(\alpha(1-\alpha)z)\,\dd\alpha
  \nonumber\\
  &\quad - \frac{7\Psi_1(0)}{48z}
  - \frac{1}{8z^2}\int_0^1\bigl[\Psi_2(\alpha(1-\alpha)z)
  - \Psi_2(0)\bigr]\,\dd\alpha\,.
  \label{eq:H-fR}
\end{align}

Combining:
\begin{align}
  H_{R,\text{bis}}(z) &=
  \frac{1}{32}\int_0^1\psi(\alpha(1-\alpha)z)\,\dd\alpha
  + \frac{1}{8z}\int_0^1\Psi_1(\alpha(1-\alpha)z)\,\dd\alpha
  \nonumber\\
  &\quad
  + \Bigl[\frac{1}{18} - \frac{7}{48}\Bigr]\frac{\Psi_1(0)}{z}
  + \Bigl[\frac{1}{3} - \frac{1}{8}\Bigr]
  \frac{1}{z^2}\int_0^1\bigl[\Psi_2(\alpha(1-\alpha)z)
  - \Psi_2(0)\bigr]\,\dd\alpha\,.
  \label{eq:H-Rbis-combined}
\end{align}

Computing the rational coefficients:
\begin{align*}
  \frac{1}{18} - \frac{7}{48}
  &= \frac{8 - 21}{144} = -\frac{13}{144}\,,\\[4pt]
  \frac{1}{3} - \frac{1}{8}
  &= \frac{8 - 3}{24} = \frac{5}{24}\,.
\end{align*}

\begin{theorem}[Scalar spectral form factor $F_2^{(0)}$]
\label{thm:F2}
\begin{equation}
  \boxed{
  \begin{aligned}
  F_2^{(0)}(z;\xi)
  &= \frac{1}{16\pi^2}\Biggl\{
    \int_0^1 \dd\alpha\,
    \Bigl[\frac{1}{32} + \frac{\xi^2}{2}
    - \xi\,\alpha(1-\alpha)\Bigr]
    \psi(\alpha(1-\alpha)z)\\
  &\quad\quad\quad\;
    + \frac{1}{8z}\int_0^1 \dd\alpha\,
      \Psi_1(\alpha(1-\alpha)z)
    - \frac{13\,\Psi_1(0)}{144z}\\
  &\quad\quad\quad\;
    + \frac{5}{24z^2}\int_0^1 \dd\alpha\,
      \bigl[\Psi_2(\alpha(1-\alpha)z) - \Psi_2(0)\bigr]
  \Biggr\}
  \end{aligned}
  }
  \label{eq:F2}
\end{equation}
\end{theorem}

\begin{proof}
From~\eqref{eq:Fi-def} and~\eqref{eq:hR}:
\begin{equation*}
  F_2^{(0)}(z;\xi) = \frac{1}{16\pi^2}\bigl[
  H_{R,\text{bis}}(z) + \xi\,H_{R\mathbf{U}}(z)
  + \xi^2\,H_{\mathbf{U}}(z)\bigr]\,.
\end{equation*}

Substituting~\eqref{eq:H-Rbis-combined}, \eqref{eq:H-RU},
and~\eqref{eq:H-U}, the $\psi$-type integrals combine:
\begin{equation*}
  \frac{1}{32}\int\psi\,\dd\alpha
  + \xi^2\cdot\frac{1}{2}\int\psi\,\dd\alpha
  - \xi\int\alpha(1-\alpha)\,\psi\,\dd\alpha
  = \int\Bigl[\frac{1}{32} + \frac{\xi^2}{2}
  - \xi\,\alpha(1-\alpha)\Bigr]\psi\,\dd\alpha\,.
\end{equation*}

The $\Psi_1$-type and $\Psi_2$-type integrals (coming only from
$H_{R,\text{bis}}$) are $\xi$-independent.  This gives~\eqref{eq:F2}.
\end{proof}


\subsection{Consistency check: exponential spectral function}

For the exponential spectral function $\psi(u) = e^{-u}$:
\begin{equation}
  \Psi_1(u) = \Psi_2(u) = e^{-u}\,,\qquad
  \Psi_1(0) = \Psi_2(0) = 1\,.
  \label{eq:exp-psi}
\end{equation}

\paragraph{$F_1^{(0)}$ cross-check:}
\begin{align*}
  F_1^{(0)}(z) &= \frac{1}{16\pi^2}\Bigl[
  \frac{1}{12z}
  + \frac{1}{2z^2}\int_0^1\bigl[e^{-\alpha(1-\alpha)z} - 1\bigr]
  \,\dd\alpha\Bigr]\\
  &= \frac{1}{16\pi^2}\Bigl[
  \frac{1}{12z} + \frac{\varphi(z)-1}{2z^2}\Bigr]
  = \frac{h_C^{(0)}(z)}{16\pi^2}\,.
\end{align*}

\paragraph{$F_2^{(0)}$ cross-check:}
Similarly, all spectral integrals reduce to the $\varphi$-based
form factors, and
\begin{equation}
  F_2^{(0)}(z;\xi)\Big|_{\psi=e^{-u}}
  = \frac{h_R^{(0)}(z;\xi)}{16\pi^2}\,.
  \label{eq:F2-exp-check}
\end{equation}

\verified{Numerical comparison of $F_1^{(0)}$ and $F_2^{(0)}$ from
the spectral integral formulas against $h_i/(16\pi^2)$ for
$\psi = e^{-u}$ at 7 values of $z \in [0.01, 100]$ and
$\xi \in \{0, 1/6, 1/4, 1\}$.  All differences $< 10^{-14}$.}


%%======================================================================
\section{Limiting Cases and UV Behaviour}
\label{sec:limits}
%%======================================================================

\subsection{Flat space limit ($R = 0$)}

When $R_{\mu\nu\rho\sigma} = 0$ (flat space), both $C^2 = 0$
and $R^2 = 0$, so the quadratic curvature
action~\eqref{eq:eff-action-weyl} vanishes identically
regardless of the form factors.  This is a trivial but mandatory
consistency check. \quad$\checkmark$


\subsection{Large-$x$ asymptotics (UV regime)}

Using $\varphi(x) = 2/x + 4/x^2 + O(x^{-3})$ as $x \to \infty$:

\paragraph{$h_C^{(0)}$:}
\begin{align}
  h_C^{(0)}(x) &= \frac{1}{12x}
  + \frac{2/x - 1 + 4/x^2 + \cdots}{2x^2}
  \nonumber\\
  &= \frac{1}{12x} - \frac{1}{2x^2} + \frac{1}{x^3} + O(x^{-4})\,.
  \label{eq:hC-UV}
\end{align}

Leading behaviour: $h_C^{(0)}(x) \sim \frac{1}{12x}$ as
$x \to \infty$.

\paragraph{$h_R^{(0)}$:}
Using the large-$x$ limits from CZ~\cite{CZ2012} eq.~(5.9):
\begin{align*}
  f_{\text{Ric}}(x) &\to \frac{1}{6x} - \frac{1}{x^2} + O(x^{-3})\,,\\
  f_R(x) &\to -\frac{1}{12x} + \frac{1}{2x^2} + O(x^{-3})\,,\\
  f_{R\mathbf{U}}(x) &\to -\frac{2}{x^2} + O(x^{-3})\,,\\
  f_{\mathbf{U}}(x) &\to \frac{1}{x} + O(x^{-2})\,.
\end{align*}

Therefore:
\begin{align}
  f_{R,\text{bis}}(x) &\to \frac{1}{18x} - \frac{1}{3x^2}
  - \frac{1}{12x} + \frac{1}{2x^2}
  = -\frac{1}{36x} + \frac{1}{6x^2} + O(x^{-3})\,,
  \label{eq:fRbis-UV}\\[4pt]
  h_R^{(0)}(x;\xi) &\to \frac{\xi^2}{x}
  - \frac{1}{36x}
  + \frac{1/6 - 2\xi + 2\xi^2}{x^2} + O(x^{-3})\,.
  \label{eq:hR-UV}
\end{align}

For general $\xi \neq 0$, the leading UV behaviour is
$h_R^{(0)} \sim \xi^2/x$, dominated by the
$f_{\mathbf{U}}$ term.  At conformal coupling $\xi = 1/6$:
\begin{equation}
  h_R^{(0)}(x;1/6) \to \frac{1/36 - 1/36}{x}
  + \frac{1/6 - 1/3 + 1/18}{x^2}
  = -\frac{1}{9x^2} + O(x^{-3})\,,
  \label{eq:hR-UV-conformal}
\end{equation}
showing \textbf{enhanced UV suppression} at conformal coupling
($1/x^2$ vs $1/x$).


\subsection{Minimal coupling ($\xi = 0$)}

With $\xi = 0$:
\begin{equation}
  h_R^{(0)}(x;0) = f_{R,\text{bis}}(x)
  = \frac{\varphi}{32}
  + \frac{\varphi/8 - 13/144}{x}
  + \frac{5(\varphi-1)}{24x^2}\,.
  \label{eq:hR-minimal}
\end{equation}
UV: $h_R^{(0)}(x;0) \to -1/(36x) + O(x^{-2})$.
Local limit: $h_R^{(0)}(0;0) = 1/72 = (1/2)(1/6)^2$.

\subsection{Conformal coupling ($\xi = 1/6$)}

At $\xi = 1/6$: $\xitilde = 0$, and the $R^2$ form factor
simplifies.  Using the explicit
forms~\eqref{eq:fRbis-explicit}--\eqref{eq:xi2-fU}:
\begin{align}
  h_R^{(0)}(x;1/6) &= f_{R,\text{bis}}(x)
  + \frac{1}{6}\,f_{R\mathbf{U}}(x)
  + \frac{1}{72}\,f_{\mathbf{U}}(x)
  \nonumber\\
  &= A(x) + \frac{1}{6}\,B(x) + \frac{1}{36}\,C(x)\,.
  \label{eq:hR-conformal}
\end{align}

We verify: at $x = 0$, this gives $1/72 - 1/36 + 1/72 = 0$.
At large $x$, this gives $-1/(6x^2) + O(x^{-3})$, confirming
enhanced UV suppression.

The conformal scalar effective action at $\mathcal{O}(\mathcal{R}^2)$
thus contains only $C^2$ at the logarithmic level (no $R^2$ divergence),
consistent with Weyl invariance.  The nonlocal $R^2$ term is
present but suppressed as $1/x^2$ in the UV.


\subsection{Dimensional analysis}

In natural units ($c = \hbar = 1$), all form factors are
\textbf{dimensionless}:
\begin{itemize}[nosep]
\item The heat kernel parameter $s$ has dimension
  $[\text{length}]^2 = [\text{mass}]^{-2}$.
\item The argument $x = s\Box$ or $x = \Box/m^2$ is dimensionless.
\item $\varphi(x)$, $f_i(x)$, $h_C^{(0)}(x)$, $h_R^{(0)}(x)$
  are all dimensionless.
\item The spectral action argument $z = \Box/\Lambda^2$
  is dimensionless.
\item $F_1^{(0)}(z)$ and $F_2^{(0)}(z)$ are dimensionless
  (checked: $[1/(16\pi^2)] \times [\psi] \times [\dd\alpha] = 1$).
\end{itemize}
\quad$\checkmark$


%%======================================================================
\section{Summary of Results}
\label{sec:summary}
%%======================================================================

\subsection{Heat kernel form factors (boxed)}

\begin{tcolorbox}[colback=blue!5!white,colframe=blue!50!black,
  title={\textbf{Scalar heat kernel form factors}}]
\textbf{Weyl sector} (independent of $\xi$):
\begin{equation}
  h_C^{(0)}(x)
  = \frac{1}{12x} + \frac{\varphi(x) - 1}{2x^2}\,,
  \qquad h_C^{(0)}(0) = \frac{1}{120}\,.
  \tag{\ref{eq:hC}}
\end{equation}

\textbf{$R^2$ sector} (depends on $\xi$):
\begin{equation}
  h_R^{(0)}(x;\xi)
  = f_{R,\text{bis}}(x)
  + \xi\,f_{R\mathbf{U}}(x) + \xi^2\,f_{\mathbf{U}}(x)\,,
  \qquad h_R^{(0)}(0;\xi) = \frac{1}{2}\Bigl(\xi-\frac{1}{6}\Bigr)^2\,.
  \tag{\ref{eq:hR}}
\end{equation}

\textbf{Component form factors}:
\begin{align*}
  f_{R,\text{bis}}(x) &= \frac{\varphi}{32}
  + \frac{\varphi/8 - 13/144}{x}
  + \frac{5(\varphi-1)}{24x^2}\,,\\[4pt]
  f_{R\mathbf{U}}(x) &= -\frac{\varphi}{4}
  - \frac{\varphi-1}{2x}\,,\\[4pt]
  f_{\mathbf{U}}(x) &= \frac{\varphi}{2}\,.
\end{align*}

\textbf{Master function}:
$\displaystyle\varphi(x) = \int_0^1 e^{-\alpha(1-\alpha)x}\,\dd\alpha$.
\end{tcolorbox}

\subsection{Spectral action form factors (boxed)}

\begin{tcolorbox}[colback=green!5!white,colframe=green!50!black,
  title={\textbf{Scalar spectral action form factors}}]
\begin{equation}
  F_1^{(0)}(z)
  = \frac{1}{16\pi^2}\Biggl\{
    \frac{\Psi_1(0)}{12z}
    + \frac{1}{2z^2}\int_0^1
      \bigl[\Psi_2(\alpha(1-\alpha)z) - \Psi_2(0)\bigr]
      \,\dd\alpha
  \Biggr\}
  \tag{\ref{eq:F1}}
\end{equation}
\begin{equation}
  \begin{aligned}
  F_2^{(0)}(z;\xi)
  &= \frac{1}{16\pi^2}\Biggl\{
    \int_0^1
    \Bigl[\frac{1}{32} + \frac{\xi^2}{2}
    - \xi\,\alpha(1-\alpha)\Bigr]
    \psi(\alpha(1-\alpha)z)\,\dd\alpha\\
  &\quad\quad\quad\;
    + \frac{1}{8z}\int_0^1
      \Psi_1(\alpha(1-\alpha)z)\,\dd\alpha
    - \frac{13\,\Psi_1(0)}{144z}\\
  &\quad\quad\quad\;
    + \frac{5}{24z^2}\int_0^1
      \bigl[\Psi_2(\alpha(1-\alpha)z) - \Psi_2(0)\bigr]
      \,\dd\alpha
  \Biggr\}
  \end{aligned}
  \tag{\ref{eq:F2}}
\end{equation}

\textbf{Spectral function definitions}:
$\psi(u) = \int_0^\infty \tilde{f}(t)\,e^{-tu}\,\dd t$,
\;$\Psi_1(u) = \int_u^\infty \psi(v)\,\dd v$,
\;$\Psi_2(u) = \int_u^\infty (v-u)\,\psi(v)\,\dd v$.

\textbf{Exponential check}: For $\psi = e^{-u}$:
$F_i^{(0)} = h_i^{(0)}/(16\pi^2)$.
\end{tcolorbox}


\subsection{BV parametric weights}

Three of the five CZ form factors admit polynomial parametric
weights
\begin{equation*}
  f_i(x) = \int_0^1 \Phi_i(\alpha)\,e^{-\alpha(1-\alpha)x}\,\dd\alpha\,,
\end{equation*}
with the following kernels:
\begin{align}
  \Phi_{\mathbf{U}}(\alpha) &= \frac{1}{2}\,,\qquad
  \Phi_{R\mathbf{U}}(\alpha) = -\alpha(1-\alpha)\,,\notag\\[4pt]
  \Phi_\Omega(\alpha) &= \frac{(1-2\alpha)^2}{4}\,.
  \label{eq:BV-weights}
\end{align}
The form factors $f_{\text{Ric}}$ and $f_R$ require subtraction
prescriptions (distributional weights) due to their $1/x$ and $1/x^2$
singularities.


\subsection{Complete benchmark table}

\begin{center}
\renewcommand{\arraystretch}{1.3}
\begin{tabular}{llccl}
\toprule
\textbf{\#} & \textbf{Quantity} & \textbf{Target} & \textbf{Result} & \textbf{Method} \\
\midrule
1 & $\tr(\Id)$ & $1$ & $1$ & Definition \\
2 & $\Omega_{\mu\nu}$ & $0$ & $0$ & Trivial bundle \\
3 & $E$ & $-\xi R - m^2$ & $-\xi R - m^2$ & Eq.~\eqref{eq:operator-data} \\
4 & $P_{\text{BV}}|_{m=0}$ & $-\xitilde R$ & $-\xitilde R$ & Eq.~\eqref{eq:xitilde} \\
5 & $h_C^{(0)}(0)$ & $1/120$ & $1/120$ & Eq.~\eqref{eq:benchmark5} \\
6 & $h_R^{(0)}(0;0)$ & $1/72$ & $1/72$ & Eq.~\eqref{eq:benchmark6} \\
7 & $h_R^{(0)}(0;1/6)$ & $0$ & $0$ & Eq.~\eqref{eq:benchmark7} \\
8 & $\beta_R(\xi)$ & $(1/2)(\xi-1/6)^2$ & $(1/2)(\xi-1/6)^2$ & Eq.~\eqref{eq:benchmark8} \\
9 & $h_C^{(0)}(x)$ & formula & Eq.~\eqref{eq:hC} & Section~\ref{sec:weyl} \\
& $h_C$ indep.\ $\xi$ & yes & yes & Remark~\ref{rem:hC-properties} \\
& Flat space & $0$ & $0$ & Section~\ref{sec:limits} \\
& UV decay & $\sim 1/x$ & confirmed & Section~\ref{sec:limits} \\
& Dim.\ analysis & OK & OK & Section~\ref{sec:limits} \\
& $F_i = h_i/(16\pi^2)$ & for $\psi=e^{-u}$ & confirmed & Section~\ref{sec:spectral} \\
\bottomrule
\end{tabular}
\end{center}


%%======================================================================
\begin{thebibliography}{20}

\bibitem{CZ2012}
A.~Codello and O.~Zanusso,
``On the non-local heat kernel expansion,''
J.\ Math.\ Phys.\ \textbf{54} (2013) 013513
[arXiv:1203.2034 [hep-th]].

\bibitem{TSR2020}
P.~de~Morais~Teixeira, I.~L.~Shapiro, and T.~G.~Ribeiro,
``One-loop effective action: nonlocal form factors and
renormalization group,''
Grav.\ Cosmol.\ \textbf{26} (2020) 185
[arXiv:2003.04503 [hep-th]].

\bibitem{BW2023}
A.~O.~Barvinsky and W.~Wachowski,
``Notes on conformal anomaly, nonlocal effective action and the
metamorphosis of the running scale,''
arXiv:2306.03780 [hep-th].

\bibitem{Vassilevich2003}
D.~V.~Vassilevich,
``Heat kernel expansion: user's manual,''
Phys.\ Rept.\ \textbf{388} (2003) 279
[hep-th/0306138].

\bibitem{BV1987}
A.~O.~Barvinsky and G.~A.~Vilkovisky,
``The generalized Schwinger--DeWitt technique in gauge theories
and quantum gravity,''
Phys.\ Rept.\ \textbf{119} (1985) 1;
Nucl.\ Phys.\ B \textbf{282} (1987) 163.

\bibitem{BV1990a}
A.~O.~Barvinsky and G.~A.~Vilkovisky,
``Covariant perturbation theory (II).\ Second order in the curvature.
General algorithms,''
Nucl.\ Phys.\ B \textbf{333} (1990) 471.

\bibitem{NT1b_conventions}
D.~Alfyorov,
``NT-1b Scalar Form Factors: Literature Conventions and Benchmarks,''
working document, March 2026.

\bibitem{NT1}
D.~Alfyorov,
``NT-1: Nonlocal Form Factors of the SCT Spectral Action ---
From Heat Kernel to Spectral Action for the Dirac Operator,''
working document, March 2026.

\end{thebibliography}

\end{document}
